% 数学ノート集・共通テンプレート v2（積分論2026年度第1部の版面を参考）
% A5判。章内で定義・定理・例・注意等の番号を共通化。
% コンパイル： platex -interaction=nonstopmode math_note_template_v2.tex (2回)
%              dvipdfmx math_note_template_v2.dvi
\documentclass[a5paper,10pt,onecolumn]{jsarticle}
\usepackage[dvipdfmx,top=22mm,bottom=20mm,inner=19mm,outer=18mm,headheight=13pt,headsep=6mm]{geometry}
\usepackage{amsmath,amssymb,amsthm,mathtools}
\usepackage{enumitem,fancyhdr,titlesec}
\usepackage[dvipdfmx,bookmarks=true,colorlinks=false]{hyperref}
\usepackage{pxjahyper}

% ----- 差し替える文書情報 -----
\newcommand{\NoteTitle}{極限・余極限と随伴函手}
\newcommand{\NoteSeries}{数学ノート集}

% ----- 文字・組版 -----
% 参考PDFの質素な明朝＋Computer Modernの印象を引き継ぎながら，
% A5判で窮屈にならないよう本文の行間を少し確保する．
\linespread{1.13}
\setlength{\parindent}{1zw}
\setlength{\parskip}{0pt}
\setlength{\abovedisplayskip}{6pt plus 2pt minus 2pt}
\setlength{\belowdisplayskip}{6pt plus 2pt minus 2pt}
\setlength{\abovedisplayshortskip}{3pt plus 2pt minus 1pt}
\setlength{\belowdisplayshortskip}{4pt plus 2pt minus 1pt}
\setlist[enumerate]{topsep=3pt,itemsep=1pt,parsep=0pt,leftmargin=2.4zw}
\titleformat{\section}{\normalfont\normalsize\bfseries}{\thesection}{0.8zw}{}
\titlespacing*{\section}{0pt}{2.0ex plus .4ex}{.9ex}
\titleformat{\subsection}{\normalfont\normalsize\bfseries}{\thesubsection}{0.7zw}{}
\titlespacing*{\subsection}{0pt}{1.5ex plus .3ex}{.6ex}
\raggedbottom

% ----- 柱とページ番号：ページ上部、罫線なし -----
\fancypagestyle{mathnote}{%
  \fancyhf{}%
  \fancyhead[L]{\footnotesize \NoteSeries\ （\NoteTitle）}%
  \fancyhead[R]{\footnotesize \thepage}%
  \renewcommand{\headrulewidth}{0pt}%
  \renewcommand{\footrulewidth}{0pt}%
}
\pagestyle{mathnote}

% ----- 節内通し番号の定義・定理群 -----
\newtheoremstyle{mathnotetheorem}
 {4.5pt}{4.5pt}{\normalfont}{}%
 {\bfseries}{}{0.65zw}%
 {\thmname{#1}\thmnumber{ #2}\thmnote{（#3）}}
\theoremstyle{mathnotetheorem}
\newtheorem{theorem}{定理}[section]
\newtheorem{definition}[theorem]{定義}
\newtheorem{lemma}[theorem]{補題}
\newtheorem{proposition}[theorem]{命題}
\newtheorem{corollary}[theorem]{系}
\newtheorem{example}[theorem]{例}
\newtheorem{remark}[theorem]{注意}
\renewcommand{\proofname}{証明}
\renewcommand{\qedsymbol}{$\square$}
\makeatletter
\renewenvironment{proof}[1][\proofname]{%
  \par\pushQED{\qed}\normalfont\topsep4\p@\@plus2\p@\relax
  \trivlist\item[\hskip\labelsep\bfseries #1]\ignorespaces
}{\popQED\endtrivlist\@endpefalse}
\makeatother

% ----- 使うノートで適宜拡張する記号 -----
\newcommand{\R}{\mathbb{R}}
\newcommand{\norm}[1]{\lVert #1\rVert}
\newcommand{\abs}[1]{\lvert #1\rvert}

\begin{document}
\section*{読了後にできること}
本稿を読むための前提は，随伴函手，普遍性とHom集合の自然同型．
本稿の具体的な到達課題は，自由函手と忘却函手の随伴を具体的な元の対応で表す．

\section{図式と錐}

本稿では圏・函手・自然変換の定義と，
米田の補題の基本事項を既知とする．
圏は特に断らない限り局所小とし，
図式の添字圏は小圏であるとする．

\begin{definition}
小圏 $J$ と圏 $\mathcal C$ に対して，
函手 $D:J\to\mathcal C$ を\emph{図式}という．
対象 $L\in\mathcal C$ から図式 $D$ への
\emph{錐}とは，各 $j\in J$ に対する射
$\lambda_j:L\to D(j)$ の族で，
任意の射 $u:j\to k$ に対して
\[
 D(u)\circ\lambda_j=\lambda_k
\]
を満たすものをいう．
\end{definition}

\begin{definition}
図式 $D$ の\emph{極限}とは，
錐 $(L,\lambda_j)$ で，
任意の錐 $(X,\xi_j)$ に対して
ただ一つの射 $f:X\to L$ が存在して
\[
 \lambda_j\circ f=\xi_j\qquad(j\in J)
\]
となるものをいう．
これを $L=\varprojlim D$ と書く．
\end{definition}

\begin{theorem}[極限の一意性]\label{thm:limit-unique}
同じ図式 $D$ の二つの極限
$(L,\lambda_j)$，$(L',\lambda'_j)$ が存在するとき，
射 $u:L\to L'$ で
$\lambda'_j u=\lambda_j$ を全ての $j$ で満たすものが
ただ一つ存在し，これは同型射である．
\end{theorem}
\begin{proof}
$L'$ の普遍性を $L$ の錐に適用すると
一意な $u:L\to L'$ が存在する．
同様に $L$ の普遍性によって
$v:L'\to L$ が存在する．
全ての $j$ について
\[
 \lambda_j vu=\lambda'_j u=\lambda_j
\]
であり，$1_L$ も同じ性質を満たす．
普遍性の一意性から $vu=1_L$ となる．
同様に $uv=1_{L'}$ だから $u$ は同型である．
\end{proof}

\section{積と等化子}

\begin{definition}
対象 $A,B$ の\emph{積}とは，
射 $p_A:P\to A$，$p_B:P\to B$ を持つ対象 $P$ で，
任意の $f:X\to A$，$g:X\to B$ に対し
唯一の $h:X\to P$ が存在して
$p_Ah=f,p_Bh=g$ となるものをいう．
これを $P=A\times B$ と書く．
\end{definition}

\begin{proposition}[積は離散図式の極限]\label{prop:product}
二つの対象からなる離散圏 $J$ と
$D:J\to\mathcal C$ に対して，
$D$ の極限は $D$ の二対象の積に他ならない．
\end{proposition}
\begin{proof}
離散圏には恒等射以外の射がないので，
錐の可換性条件は空であり，
錐は二本の射 $X\to A$，$X\to B$ と同じである．
極限の普遍性をそのまま書けば
積の普遍性に一致する．
\end{proof}

\begin{definition}
平行な二射 $f,g:A\to B$ の\emph{等化子}とは，
射 $e:E\to A$ で $fe=ge$ を満たし，
$fh=gh$ となる任意の $h:X\to A$ に対して
一意な $\bar h:X\to E$ が存在して
$e\bar h=h$ となるものをいう．
\end{definition}

\begin{proposition}[集合の等化子]\label{prop:equalizer-set}
$\mathbf{Set}$ において
\[
 E=\{a\in A:f(a)=g(a)\}
\]
を取り，包含写像 $e:E\hookrightarrow A$ とすると
$(E,e)$ は $f,g$ の等化子である．
\end{proposition}
\begin{proof}
定義から $fe=ge$ である．
$h:X\to A$ が $fh=gh$ を満たせば，
任意の $x\in X$ について
$f(h(x))=g(h(x))$ だから $h(x)\in E$ である．
従って $\bar h(x)=h(x)$ と置けば
$X\to E$ の写像が得られ $e\bar h=h$ となる．
$e$ は単射なので $\bar h$ は一意である．
\end{proof}

\begin{theorem}[集合の小図式の極限]\label{thm:set-limits}
小圏 $J$ と函手 $D:J\to\mathbf{Set}$ に対して，
極限 $\varprojlim D$ は常に存在する．
\end{theorem}
\begin{proof}
積集合
\[
 P=\prod_{j\in\operatorname{Ob}(J)}D(j)
\]
を取る．$J$ は小圏なので対象全体は集合であり，
選択公理を用いて積を構成できる．
\[
 L=\{(x_j)_j\in P:
 D(u)(x_j)=x_k
 \text{ for all }u:j\to k\}
\]
と置く．$L$ から $D(j)$ への射を座標射影 $\lambda_j$ とする．
$L$ の定義から
$D(u)\lambda_j=\lambda_k$ だから錐である．
任意の錐 $(X,\xi_j)$ に対し
\[
 h:X\to L,\qquad h(x)=(\xi_j(x))_j
\]
と定めることができる．
実際，錐の可換性から右辺は $L$ の元となる．
$\lambda_jh=\xi_j$ が成り立つ．
この条件を満たす写像は各座標の値によって一意なので，
$h$ も一意である．従って $L$ は極限である．
\end{proof}

\section{余極限と押出し}

\begin{definition}
図式 $D:J\to\mathcal C$ から対象 $Q$ への
\emph{余錐}とは，射 $\iota_j:D(j)\to Q$ の族で
任意の $u:j\to k$ について
$\iota_kD(u)=\iota_j$ を満たすものをいう．
任意の余錐 $(Q',\iota'_j)$ に対して
唯一の $h:Q\to Q'$ が存在して
$h\iota_j=\iota'_j$ となる余錐 $(Q,\iota_j)$ を
\emph{余極限}といい，
$Q=\varinjlim D$ と書く．
\end{definition}

\begin{theorem}[集合の小図式の余極限]\label{thm:set-colimits}
小圏 $J$ と函手 $D:J\to\mathbf{Set}$ に対し，
余極限 $\varinjlim D$ は常に存在する．
\end{theorem}
\begin{proof}
非交和
$S=\bigsqcup_{j\in\operatorname{Ob}(J)}D(j)$ を取り，
$D(j)$ の元 $x$ を $(j,x)$ と表す．
$J$ の任意の射 $u:j\to k$ と $x\in D(j)$ に対して
\[
 (j,x)\sim_0(k,D(u)(x))
\]
という関係を考え，
その反射・対称・推移閉包を $\sim$ とする．
商集合 $Q=S/{\sim}$ とし，
$\iota_j(x)=[(j,x)]$ と置く．
定義から $\iota_kD(u)=\iota_j$ だから余錐である．

別の余錐 $(Q',\iota'_j)$ を取る．
$S$ 上の写像を $h_0(j,x)=\iota'_j(x)$ と定める．
余錐の条件から
\[
 h_0(j,x)=\iota'_j(x)
 =\iota'_k(D(u)(x))
 =h_0(k,D(u)(x))
\]
だから $h_0$ は $\sim_0$ で関係づけられた点を
同じ値へ写す．
従って $\sim$ に関しても同じ値を持ち，
一意な写像 $h:Q\to Q'$ を誘導する．
$h\iota_j=\iota'_j$ となり，
$\iota_j$ の像が $Q$ を尽くすので $h$ は一意である．
\end{proof}

\begin{definition}
射 $f:A\to B$，$g:A\to C$ に対し，
$B$ と $C$ からの射を持ち
\[
 \begin{array}{ccc}
 A&\xrightarrow{f}&B\\
 {\scriptstyle g}\downarrow&&\downarrow\\
 C&\longrightarrow&P
 \end{array}
\]
が可換になる余極限を\emph{押出し}と呼ぶ．
$\mathbf{Set}$ では
$P=(B\sqcup C)/{\sim}$，
ただし $f(a)\sim g(a)$ で生成される同値関係である．
\end{definition}

\section{随伴函手}

\begin{definition}
$F:\mathcal C\to\mathcal D$，
$G:\mathcal D\to\mathcal C$ を函手とする．
各 $A\in\mathcal C$，$B\in\mathcal D$ に関して
自然な全単射
\[
 \Phi_{A,B}:
 \operatorname{Hom}_{\mathcal D}(F(A),B)
 \longrightarrow
 \operatorname{Hom}_{\mathcal C}(A,G(B))
\]
が存在するとき，$F$ は $G$ の\emph{左随伴}，
$G$ は $F$ の\emph{右随伴}であるといい，
$F\dashv G$ と書く．
\end{definition}

\begin{example}[自由群の随伴性]\label{ex:free-group}
集合から自由群を作る函手
$F:\mathbf{Set}\to\mathbf{Grp}$ と
忘却函手 $U:\mathbf{Grp}\to\mathbf{Set}$ に対して
\[
 F\dashv U
\]
が成り立つ．
\end{example}
\begin{proof}
集合 $S$ に対する自由群 $F(S)$ の普遍性は，
任意の群 $G$ に対する全単射
\[
 \operatorname{Hom}_{\mathbf{Grp}}(F(S),G)
 \cong\operatorname{Hom}_{\mathbf{Set}}(S,U(G))
\]
を与える．
左から右への写像は
自由生成元への包含 $S\to F(S)$ との合成である．
右から左への写像は自由群の普遍性が与える
準同型への一意な延長である．
集合の写像 $S'\to S$ または群準同型 $G\to G'$ と
合成しても同じ対応が得られるので，
この全単射は両変数について自然である．
従って随伴性が成立する．
\end{proof}

\begin{example}[対角函手と積]
積が存在する圏 $\mathcal C$ において，
対角函手 $\Delta:\mathcal C\to\mathcal C\times\mathcal C$ と
二項積函手
$\times:\mathcal C\times\mathcal C\to\mathcal C$
に対して
\[
 \Delta\dashv\times
\]
が成り立つ．
\end{example}
\begin{proof}
二項積の普遍性によって
\[
 \operatorname{Hom}_{\mathcal C\times\mathcal C}
 ((X,X),(A,B))
 \cong
 \operatorname{Hom}_{\mathcal C}(X,A\times B)
\]
が成立する．
左辺は
$\operatorname{Hom}(X,A)\times\operatorname{Hom}(X,B)$
なので，対応は射の対 $(f,g)$ を
その積への一意な誘導射 $\langle f,g\rangle$ へ写す．
積の普遍性から合成に関して自然であるため，
随伴性を得る．
\end{proof}

\section{単位・余単位と三角恒等式}

\begin{theorem}[随伴の単位と余単位]\label{thm:unit-counit}
$F\dashv G$ とし，
$\Phi_{A,B}$ を随伴の自然な全単射とする．
\[
 \eta_A=\Phi_{A,F(A)}(1_{F(A)}):
 A\longrightarrow GF(A)
\]
および
\[
 \varepsilon_B=\Phi^{-1}_{G(B),B}(1_{G(B)}):
 FG(B)\longrightarrow B
\]
と置くと，
$\eta:1_{\mathcal C}\Rightarrow GF$ と
$\varepsilon:FG\Rightarrow1_{\mathcal D}$
は自然変換となり，
\[
 \varepsilon_{F(A)}\circ F(\eta_A)=1_{F(A)},
 \qquad
 G(\varepsilon_B)\circ\eta_{G(B)}=1_{G(B)}
 \tag{T}
\]
を満たす．
逆にこのような自然変換 $\eta,\varepsilon$ が
(T) を満たせば，$F\dashv G$ である．
\end{theorem}
\begin{proof}
まず随伴の全単射の自然性を用いて，
任意の $f:F(A)\to B$ に対し
\[
 \Phi_{A,B}(f)=G(f)\circ\eta_A
 \tag{1}
\]
を得る．
実際 $B$ 変数に関する自然性から
\[
 \Phi_{A,B}(f)
 =G(f)\circ\Phi_{A,F(A)}(1_{F(A)}).
\]
同様に逆写像について
\[
 \Phi^{-1}_{A,B}(g)=\varepsilon_B\circ F(g)
 \qquad(g:A\to G(B))
 \tag{2}
\]
が成り立つ．

$\eta$ の自然性を示す．
$u:A\to A'$ について，
随伴の $A$ 変数に関する自然性から
\[
 GF(u)\circ\eta_A
 =\Phi_{A,F(A')}(F(u))
 =\eta_{A'}\circ u
\]
となる．
$\varepsilon$ の自然性も同じ議論を逆写像に適用して
示される．
式 (2) を $g=\eta_A$，$B=F(A)$ に適用すると
\[
 1_{F(A)}
 =\Phi^{-1}_{A,F(A)}(\eta_A)
 =\varepsilon_{F(A)}\circ F(\eta_A)
\]
を得る．
同様に式 (1) を
$f=\varepsilon_B$，$A=G(B)$ に適用すると
\[
 1_{G(B)}
 =\Phi_{G(B),B}(\varepsilon_B)
 =G(\varepsilon_B)\circ\eta_{G(B)}
\]
を得る．これが (T) である．

逆に $\eta,\varepsilon$ が与えられたとする．
任意の $A,B$ に対し
\[
 \Phi(f)=G(f)\circ\eta_A,\qquad
 \Psi(g)=\varepsilon_B\circ F(g)
\]
と定める．
自然性を用いると
\begin{align*}
 \Psi(\Phi(f))
 &=\varepsilon_B F(G(f)\eta_A)\\
 &=\varepsilon_B FG(f)F(\eta_A)\\
 &=f\varepsilon_{F(A)}F(\eta_A)=f,
\end{align*}
ただし三行目で $\varepsilon$ の自然性を用いた．
同様に
\begin{align*}
 \Phi(\Psi(g))
 &=G(\varepsilon_B)GF(g)\eta_A\\
 &=G(\varepsilon_B)\eta_{G(B)}g=g
\end{align*}
となる．最後の二行目では $\eta$ の自然性を用いた．
従って $\Phi$ と $\Psi$ は互いに逆な全単射である．
定義から両変数に関する自然性も従うので
$F\dashv G$ を得る．
\end{proof}

\section{随伴と極限・余極限の保存}

\begin{theorem}[左随伴は余極限を保存する]
\label{thm:left-preserves}
$F:\mathcal C\to\mathcal D$ が
$G:\mathcal D\to\mathcal C$ の左随伴であるとする．
小図式 $D:J\to\mathcal C$ の余極限
$(L,\iota_j)$ が存在するとき，
$(F(L),F(\iota_j))$ は図式 $F\circ D$ の余極限となる．
\end{theorem}
\begin{proof}
$B\in\mathcal D$ とし，
$(B,\beta_j:F(D(j))\to B)$ を
$F\circ D$ の余錐とする．
随伴の全単射によって各 $\beta_j$ は
射
\[
 \gamma_j=\Phi_{D(j),B}(\beta_j):
 D(j)\longrightarrow G(B)
\]
に一意に対応する．
$\beta_kF(D(u))=\beta_j$ という余錐条件は，
随伴の自然性によって
$\gamma_kD(u)=\gamma_j$ と同値である．
従って $(G(B),\gamma_j)$ は $D$ の余錐となる．
$L$ の普遍性から一意な射 $g:L\to G(B)$ が存在して
$g\iota_j=\gamma_j$ となる．
これを随伴の逆全単射で戻すと
\[
 f=\Phi^{-1}_{L,B}(g):F(L)\longrightarrow B
\]
が得られる．
自然性から
$fF(\iota_j)=\beta_j$ である．
逆の対応も全単射なので，この $f$ は一意である．
従って $(F(L),F(\iota_j))$ は余極限である．
\end{proof}

\begin{theorem}[右随伴は極限を保存する]
\label{thm:right-preserves}
$F\dashv G$ とし，
小図式 $D:J\to\mathcal D$ の極限
$(L,\lambda_j)$ が存在するとする．
このとき
$(G(L),G(\lambda_j))$ は
$G\circ D$ の極限である．
\end{theorem}
\begin{proof}
任意の $X\in\mathcal C$ と
錐 $(X,\xi_j:X\to G(D(j)))$ を取る．
随伴の逆全単射により各 $\xi_j$ は
射
$\beta_j:F(X)\to D(j)$ に一意に対応する．
自然性によって錐の可換性は
$D(u)\beta_j=\beta_k$ と同値である．
従って $\beta_j$ は $D$ への錐となる．
$L$ の極限の普遍性から一意な
$h:F(X)\to L$ で
$\lambda_jh=\beta_j$ となるものが存在する．
$h$ を随伴で移せば一意な
$\bar h:X\to G(L)$ が得られ，
$G(\lambda_j)\bar h=\xi_j$ となる．
よって $G(L)$ は極限を表す．
\end{proof}

\begin{remark}
極限・余極限は「図式からの錐」または
「図式への余錐」を普遍的に表す対象である．
随伴はHom集合の自然同型によって射の向きを
二つの圏の間で移す操作であり，
この自然性によって極限・余極限の保存則が従う．
\end{remark}

\section{到達課題の完全解答}
\noindent\textbf{自由群の随伴を具体化する．}
\begin{proposition}
集合 $S$ と群 $G$ について
\[
\operatorname{Hom}_{\mathbf{Grp}}(F(S),G)
\cong
\operatorname{Hom}_{\mathbf{Set}}(S,U(G))
\]
が自然に成立する．ここで $F(S)$ は自由群，
$U$ は忘却函手である．
\end{proposition}
\begin{proof}
左辺の準同型 $\varphi$ を生成元集合に制限し，
$s\mapsto\varphi(s)$ とする．
逆に写像 $f:S\to G$ が与えられれば，
自由群の普遍性からただ一つの準同型
$\widetilde f:F(S)\to G$ が存在する．
制限と延長は互いに逆である．
集合写像と群準同型の合成に関してこの対応は可換なので
両変数で自然である．
\end{proof}
\noindent\textbf{到達確認．}
$S=\{s\}$ のとき $F(S)$ の同型類を求めよ．
\par\noindent\textbf{解答．}
一生成元の自由群は無限巡回群
$\mathbb Z$ に同型である．任意の群 $G$ への
準同型は1の像を任意に選ぶことで一意に決まる．
\end{document}